\section{Heat Transfer \& Thermal Protection}

\subsection{Introduction}
This chapter details the heat transfer mechanisms and thermal protection systems (TPS) used in the rocket engine. The models cover convective and radiative heat transfer from the hot gas, and various cooling methods including film, regenerative, and ablative cooling.

\subsection{Nomenclature}
\begin{description}
    \item[\(q''\)] Heat flux [\si{W/m^2}]
    \item[\(h_g\)] Gas-side heat transfer coefficient [\si{W/(m^2.K)}]
    \item[\(T_{aw}\)] Adiabatic wall temperature [\si{K}]
    \item[\(T_w\)] Wall temperature [\si{K}]
    \item[\(\epsilon\)] Emissivity [dimensionless]
    \item[\(\sigma\)] Stefan-Boltzmann constant [\si{W/(m^2.K^4)}]
    \item[\(B\)] Blowing parameter [dimensionless]
    \item[\(\dot{r}\)] Recession rate [\si{m/s}]
    \item[\(H_{abl}\)] Heat of ablation [\si{J/kg}]
\end{description}

\subsection{Gas-Side Heat Transfer}

\subsubsection{Convective Heat Flux}
The convective heat flux from the gas to the wall is:
\begin{equation}
q''_{conv} = h_g (T_{aw} - T_w)
\end{equation}
\coderef{engine/pipeline/thermal/ablative_cooling.py}{32}

\subsubsection{Adiabatic Wall Temperature}
The adiabatic wall temperature accounts for kinetic energy recovery:
\begin{equation}
T_{aw} = T_c \times r_{rec}
\end{equation}
where \(r_{rec}\) is the recovery factor.
\coderef{engine/pipeline/thermal/regen_cooling.py}{459, 601}

\subsubsection{Radiative Heat Flux}
The radiative heat flux from the gas to the wall is:
\begin{equation}
q''_{rad} = \epsilon \sigma (T_{gas}^4 - T_w^4)
\end{equation}
\coderef{engine/pipeline/thermal/ablative_cooling.py}{37}

\subsection{Cooling Models}

\subsubsection{Film Cooling}
Film cooling effectiveness \(\eta_{film}\) reduces the effective gas temperature:
\begin{equation}
T_{gas,eff} = T_c - \eta_{film} (T_c - T_{inj,film})
\end{equation}
\coderef{engine/pipeline/thermal/film_cooling.py}{124}

\subsubsection{Regenerative Cooling}
Regenerative cooling uses a thermal resistance network:
\begin{equation}
U = \left( \frac{1}{h_g} + \frac{t_{wall}}{k_{wall}} + \frac{1}{h_{coolant}} \right)^{-1}
\end{equation}
\coderef{engine/pipeline/thermal/regen_cooling.py}{485-486}

\subsubsection{Ablative Cooling}
Ablation occurs when \(T_w > T_{pyro}\). The mass flux of ablated material is:
\begin{equation}
\dot{m}''_{abl} = \frac{q''_{eff}}{H_{abl} + c_p (T_w - T_{pyro})}
\end{equation}
The recession rate is \(\dot{r} = \dot{m}''_{abl} / \rho_{mat}\).
\coderef{engine/pipeline/thermal/ablative_cooling.py}{225-237}

\subsubsection{Blowing Effect}
Pyrolysis gases reduce convective heat transfer via the blowing parameter \(B\):
\begin{equation}
q''_{conv,eff} = q''_{conv} \times \frac{1}{1 + c_{blowing} B}
\end{equation}
where \(B = \dot{m}_{pyro} / \dot{m}_{ext}\).
\coderef{engine/pipeline/thermal/ablative_cooling.py}{188-192}

\subsection{Graphite Oxidation}
Graphite recession is primarily driven by oxidation. The rate is the minimum of kinetic and diffusion limits:
\begin{equation}
\dot{m}''_{ox} = \min(\dot{m}''_{ox,kin}, \dot{m}''_{ox,diff})
\end{equation}
\coderef{engine/pipeline/thermal/graphite_cooling.py}{333}

The kinetic rate follows an Arrhenius form:
\begin{equation}
\dot{m}''_{ox,kin} = j_{ref} \exp\left( -\frac{E_a}{R_{gas}} \left( \frac{1}{T_w} - \frac{1}{T_{ref}} \right) \right) \left( \frac{p_{O2}}{P_{ref}} \right)^n
\end{equation}
\coderef{engine/pipeline/thermal/graphite_cooling.py}{297}

